\documentclass[11pt]{article}
\usepackage[margin=1in]{geometry}
\usepackage[T1]{fontenc}
\usepackage{lmodern,amsmath,amssymb,amsthm,booktabs,microtype,xcolor}
\usepackage[colorlinks=true,urlcolor=blue!50!black,linkcolor=black]{hyperref}
\newtheorem{theorem}{Theorem}
\newtheorem{lemma}[theorem]{Lemma}
\newcommand{\repo}{https://github.com/Akhil-rav/discrete-math-lab}
\hypersetup{pdftitle={Fifteen extremal colorings of the ladder}}
\title{\LARGE\bfseries Fifteen extremal colorings of the ladder}
\author{}
\date{Working note \(\,\cdot\,\) November 16, 2026}
\begin{document}
\maketitle
\begin{abstract}
How many ways can a ladder be colored with as many colors as possible while
every perfect matching repeats a color? From five columns onward there are
exactly fifteen edge partitions. All repeated colors occur at the ends.
The proof contracts an interior square and checks a finite list of base cases
by exhaustive computation. A separate argument bounds the number of horizontal
tiles needed to witness a rainbow matching.
\end{abstract}

\section{The fifteen partitions}
Write \(L_k=P_k\square K_2\) and put \(s=k-1\). Label the rungs
\(r_0,\ldots,r_s\), and let \(H_i\) be the pair of horizontal edges between
columns \(i\) and \(i+1\). A perfect matching tiles the columns by single
rungs and disjoint horizontal pairs. The leftmost free column forces one of
these two choices, after which the same argument applies to the remaining
columns. A matching is rainbow if its edges have distinct colors.

\begin{theorem}
For every \(k\ge5\), the following list gives exactly the colorings with
\(3k-5\) colors and no rainbow perfect matching, up to renaming colors.
Edges are fixed. There are fifteen partitions and six orbits under graph
automorphisms.
\end{theorem}
Each table entry lists the nonsingleton color classes. Every other edge
gets its own color. Reversal interchanges the ends of the ladder.
\begin{center}\small
\begin{tabular}{@{}lll@{}}
\toprule
Classes & Choices & Count\\
\midrule
\(\{r_0,r_s,x,y\}\) & \(x\in H_0,\ y\in H_{s-1}\) & 4\\
\(\{r_0,r_1,x\},\,H_0\) & \(x\in H_1\), also reverse & 4\\
\(\{r_0,r_s,x\},\,H_{s-1}\) & \(x\in H_0\), also reverse & 4\\
\(\{r_0,r_1\},\,H_0,\,H_1\) & also reverse & 2\\
\(\{r_0,r_s\},\,H_0,\,H_{s-1}\) & & 1\\
\bottomrule
\end{tabular}
\end{center}
For the first row, a matching takes either \(r_0\) or \(H_0\), and either
\(r_s\) or \(H_{s-1}\). Thus it uses two edges of the quadruple. In rows two
and four, a matching omitting \(H_0\) uses \(r_0\), followed by either
\(r_1\) or \(H_1\). In row three, omitting \(H_{s-1}\) forces \(r_s\) and
the left end supplies another edge of the triple. In the last row, omitting
both end rail pairs forces both end rungs. Every listed partition works.

\newpage
\section{Why the list is exhaustive}
The \href{\repo/blob/main/notes/matching_covers.md}{ladder bound} gives the
maximum \(3k-5\) for \(k\ge3\). Since the ladder has \(3k-2\) edges,
an extremizer has deficit three: the sum of \(|B|-1\) over its color
classes \(B\) is three. The possible nonsingleton sizes are \(4\), \(3+2\),
and \(2+2+2\). Call the edges in those classes marked. There are at most six.

\begin{lemma}
If the square between columns \(q,q+1\) has four unmarked edges, its two
columns can be merged into one, preserving all marked classes. If the
original coloring avoids rainbow perfect matchings, so does the shorter one.
\end{lemma}
\begin{proof}
Delete gap \(q\), merge its two rungs into a single fresh singleton rung,
and keep all other edges in order on their respective rails or rung line.
The shortened coloring has the same deficit. A tiling is specified by its
binary word on gaps: a one selects both rails, and consecutive ones are
forbidden. In any shortened tiling insert a zero at position \(q\).
This gives a legal original tiling. Outside the square, rail choices agree
and each rung sees the same incident gap choices. Inside the square only
singleton colors can enter. Hence a rainbow tiling on the shorter ladder
would give a rainbow tiling on the original.
\end{proof}
For \(k\ge18\), consider the \(k-5\) gaps \(2\le q\le k-4\).
Each marked edge touches at most two of their squares, so at most twelve
squares contain a marked edge. There is a square to which the lemma applies.
After contraction, apply the classification at \(k-1\). The first two gaps,
last two gaps, first two rungs, and last two rungs of the shorter ladder
correspond to those of the original. All marked edges survived. The shorter
table therefore determines the original partition uniquely for this gap.

\paragraph{Finite base.}
For each \(5\le k\le17\) the
\href{\repo/blob/main/experiments/ladder_stability.py}{exact enumeration}
compares the full set of surviving partitions with the table. It lists the
matchings by choosing all subsets of nonconsecutive gaps. For a quadruple
it checks every four-edge subset. For a triple it intersects all matchings
containing at most one of its edges. Any completing pair must lie in that
intersection and outside the triple, and every such pair works. For three
pairs, one pair \(P\) must lie in the all-rung matching. Choose any matching
that omits at least one edge of \(P\). Some second pair \(Q\) must lie in it.
The third pair is found by intersecting all matchings containing neither
whole pair. Disjointness is required throughout, and duplicate partitions
are removed. These loops exhaust the three shapes. At every length the
three counts are \(4,8,3\), with exact agreement of edge positions.
This supplies the base of the induction.

Choosing the image of one corner forces its neighbors by degree, then
forces successive columns. The four automorphisms are rail exchange, end
reversal, their product and the identity. Row one splits by whether the
selected rails lie on the same side. Each other row is one orbit, giving six.

\newpage
\section{How far away is a rainbow matching?}
The finite check used every tiling. There is also a uniform way to see why
few horizontal pairs suffice when the deficit is small.

\begin{theorem}
Let a finite simple graph have a fixed perfect matching \(R\) and an edge
coloring using \(|E|-d\) colors. If it has a rainbow
perfect matching, it has one obtained from \(R\) by flipping at most \(d\)
vertex-disjoint alternating cycles.
\end{theorem}
\begin{proof}
Take a rainbow perfect matching \(M\). In each color class keep
the edge used by \(M\), if present, and otherwise keep any one edge. Call
the remaining \(d\) edges forbidden. All kept edges have distinct colors.
The symmetric difference \(R\mathbin{\triangle}M\) splits into disjoint
even cycles alternating between the two matchings: every vertex incident
to the difference has one edge of each matching. Flip precisely the cycles
containing a forbidden edge of \(R\). Every forbidden edge of \(R\) lies on
one of them because \(M\) avoids it. At most \(d\) cycles are flipped.
All newly used edges belong to \(M\), and all retained edges avoid the
forbidden set. The resulting matching is perfect and rainbow.
\end{proof}

On the ladder take \(R\) to be all rungs. Each horizontal tile in \(M\)
gives exactly one alternating square of the symmetric difference.
Consequently, if a deficit-three coloring admits a rainbow matching,
there is one with at most three horizontal tiles. Two do not always
suffice. On eight columns, give the rung pairs
\[
 \{r_0,r_1\},\quad \{r_3,r_4\},\quad \{r_6,r_7\}
\]
three different repeated colors, with fresh colors elsewhere. A rainbow
matching must omit a rung from each pair. A horizontal tile can hit at
most one pair, and tiles across gaps \(0,3,6\) achieve the bound.

\paragraph{Reproduction.}
From the repository root run
\begin{center}\texttt{python -m experiments.ladder\_stability}\end{center}
The thirteen base cases range from 8 to 2,584 perfect matchings.
A \href{\repo/blob/main/tests/test_ladder_stability.py}{second search}
generates matchings by pairing vertices. Starting with distinct colors, it
merges two colors present in a surviving rainbow matching at each step.
Every avoiding partition can be reached this way: at each step that
matching contains two colors that coincide in the final partition. At the
last step the merged colors must both occur in every survivor. This check
returns exactly the same partitions throughout the base range.
The earlier \href{\repo/blob/main/notes/l5_extremizers.md}{five-column proof}
also derives that case by hand.

\paragraph{Related work.}
Haas and Young study a
\href{https://doi.org/10.1016/j.disc.2011.10.017}{complete-graph perfect-matching threshold}.
The published abstract of
\href{https://doi.org/10.1016/j.dam.2023.11.049}{\emph{Anti-Ramsey number of matchings in outerplanar graphs}}
reports an extremum over a family of Hamiltonian bipartite outerplanar hosts.
The classification here fixes one host at each length and retains edge
positions. Prior coverage of this exact classification remains unverified.
\end{document}
